\section{Methodology}
\label{sec:methodology}

\begin{figure*}[t]
\centering
\begin{tikzpicture}[node distance=1.5cm, auto, >=latex, thick, scale=0.85, every node/.style={transform shape}]
    % Nodes
    \node [block] (input) {Input Feature $z(x)_j$};
    \node [block, right=0.5cm of input] (proj) {Sphere Projection $\psi(x)_j$};
    \node [block, right=0.5cm of proj] (init) {Lattice Init $s_j^{(0)}$};
    
    % Loop for layers
    \node [lattice, below=1cm of init] (state_prev) {Prev State $s_j^{(l-1)}$};
    \node [func, left=0.6cm of state_prev] (logistic) {Logistic Map $f(s)$};
    \node [block, left=0.6cm of logistic] (coupling) {Spatial Coupling $\bar{s}_j^{(l)}$};
    \node [sum, below=0.8cm of coupling] (steer) {$+$};
    \node [block, right=0.6cm of steer] (cand) {Candidate State $s_{cand}^{(l)}$};
    
    % Gated mix
    \node [sum, right=0.8cm of cand] (gate_mix) {$\oplus$};
    \node [lattice, right=0.8cm of gate_mix] (state_curr) {Curr State $s_j^{(l)}$};
    \node [block, below=1cm of state_curr] (weight) {Weight Assembly $W_{merged, j}^{(l)}$};

    % Connections
    \draw [->] (input) -- (proj);
    \draw [->] (proj) -- (init);
    \draw [->] (init) -- (state_prev);
    
    \draw [->] (state_prev) -- (logistic);
    \draw [->] (logistic) -- (coupling);
    \draw [->] (coupling) -- (steer);
    \draw [->] (proj) |- node[pos=0.7, left] {Perturbation Steering} (steer);
    \draw [->] (steer) -- (cand);
    
    \draw [->] (cand) -- (gate_mix);
    \draw [->] (state_prev) -| node[pos=0.3, above] {Skip Connection $(1-\lambda_l)$} (gate_mix);
    \draw [->] (gate_mix) -- (state_curr);
    
    \draw [->] (state_curr) -- (weight);
    
    % Feed back loop
    \draw [->] (state_curr) -- ++(0,1.2) -| (state_prev);

    % Legend or box
    \node [draw, dashed, gray, fit=(logistic) (coupling) (steer) (cand) (gate_mix), inner sep=0.25cm, label=above:{\textbf{Gated Coupled Map Lattice (G-CML) Layer $l$}}] (g_cml_box) {};
\end{tikzpicture}
\caption{The architectural pipeline of \textbf{ChaosMerge}. Spherically projected input features steer the Coupled Map Lattice (CML) trajectory. Within each G-CML layer group, states undergo non-linear spatial-chaotic coupling, local perturbation steering, and learned gating skip-connections, before scaling task-specific vectors for final model assembly.}
\label{fig:architecture}
\end{figure*}

In this section, we formulate our proposed framework, \textbf{ChaosMerge} (Chaos-Theoretic Attractor Merging). 
We treat the sequence of $L$ layer groups in a deep neural network as discrete steps in time. 
Instead of computing static merging coefficients or relying on a standard feed-forward network to route features, we model the merging coefficients as the trajectory of a non-linear, chaotic Coupled Map Lattice (CML) driven by a Logistic Map, stabilized via a novel gating mechanism to prevent gradient explosion.

Let $W_{base}^{(l)}$ represent the weights of the pre-trained base model at layer group $l \in \{1, \dots, L\}$. 
For a set of $K$ fine-tuned expert models, let $W_k^{(l)}$ be the weights of expert $k \in \{1, \dots, K\}$ at layer group $l$. 
We define the task vector $V_k^{(l)}$ for expert $k$ at layer $l$ as:
\begin{equation}
V_k^{(l)} = W_k^{(l)} - W_{base}^{(l)}.
\end{equation}
For a given input task feature, we construct a discrete lattice of $K$ nodes. 
At each layer group $l$, the state of the lattice for task $j \in \{1, \dots, K\}$ is denoted by $s_j^{(l)} \in [0, 1]^K$, where $s_{k, j}^{(l)}$ represents the state of node $k$ (corresponding to the scaling coefficient or activation level of task $k$'s expert parameters).

\subsection{Sphere-Projected Feature Extraction}
To guide the chaotic trajectory of the lattice states without introducing high-dimensional parameter inflation, we extract low-dimensional phase-space representations from the input. Let $z(x)_j \in \mathbb{R}^D$ be the spatially averaged patch tokens from the base model's frozen patch embedding layer for task $j$. We project $z(x)_j$ into a $d$-dimensional phase-space (where $d = K = 4$) via a frozen random projection matrix $P \in \mathbb{R}^{D \times d}$ and normalize it to the unit sphere:
\begin{align}
\tilde{\psi}(x)_j &= z(x)_j P \in \mathbb{R}^d, \\
\psi(x)_j &= \frac{\tilde{\psi}(x)_j}{\|\tilde{\psi}(x)_j\|_2 + \epsilon} \in \mathbb{R}^d,
\end{align}
where $\epsilon = 10^{-8}$ is a numerical stabilizer. Projecting features onto the unit sphere $\mathbb{S}^{d-1}$ serves as a powerful geometric regularizer, ensuring that input features are scale-invariant and bounded before driving the chaotic system.

\subsection{Lattice State Initialization}
At the initial temporal step $l = 0$, we initialize the lattice states $s_j^{(0)} \in [0, 1]^K$ via a learned linear projection of the input phase state:
\begin{align}
\tilde{s}_{k, j}^{(0)} &= \sum_{m=1}^d W_{init, k, m} \psi(x)_{j, m} + b_{init, k}, \\
s_{k, j}^{(0)} &= \sigma\left(\tilde{s}_{k, j}^{(0)}\right) = \frac{1}{1 + e^{-\tilde{s}_{k, j}^{(0)}}},
\end{align}
where $W_{init} \in \mathbb{R}^{K \times d}$ and $b_{init} \in \mathbb{R}^K$ are trainable initializers, and $\sigma(\cdot)$ is the standard Sigmoid activation function ensuring that initial states are strictly bounded within the interval $[0, 1]$.

\subsection{Gated Coupled Map Lattice (G-CML)}
A major technical challenge in deep recurrent lattices is gradient explosion. In a standard Coupled Map Lattice, propagating states recursively through $L=14$ layers using a chaotic Logistic Map $f(u) = 4u(1-u)$ whose derivative is bounded by $4.0$ leads to an exponential gradient multiplier scaling up to $4^{14} \approx 2.68 \times 10^8$. This inherent instability makes the network un-optimizable by first-order gradient descent.

To resolve this fundamental flaw, we introduce the \textbf{Gated Coupled Map Lattice (G-CML)}, which incorporates a learned layer-wise gating coefficient $\lambda_l \in [0, 1]$ acting as a residual skip-connection.
For each subsequent layer group $l \in \{1, \dots, L\}$, we first compute the spatial-chaotic coupling across the nodes:
\begin{equation}
\bar{s}_{k, j}^{(l)} = (1 - \gamma_l) f\left(s_{k, j}^{(l-1)}\right) + \frac{\gamma_l}{K} \sum_{i=1}^K f\left(s_{i, j}^{(l-1)}\right),
\end{equation}
where:
\begin{itemize}
    \item $f(u) = 4u(1-u)$ is the fully chaotic Logistic Map.
    \item $\gamma_l \in [0, 1]$ is a learned layer-wise spatial coupling coefficient representing the local diffusion of task nodes, modeled as $\gamma_l = \sigma(\gamma_{raw, l})$ with a raw parameter $\gamma_{raw, l} \in \mathbb{R}$.
\end{itemize}

Next, we project the coupled state into logit-space, inject the localized input perturbation to steer the trajectory, and apply the learned gating:
\begin{align}
s_{cand, k, j}^{(l)} &= \sigma\left( \sigma^{-1}\left( \text{clip}\left(\bar{s}_{k, j}^{(l)}, \delta, 1-\delta\right) \right) + \langle \psi(x)_j, \Phi_k^{(l)} \rangle + \phi_k^{(l)} \right), \\
s_{k, j}^{(l)} &= (1 - \lambda_l) s_{k, j}^{(l-1)} + \lambda_l s_{cand, k, j}^{(l)},
\end{align}
where:
\begin{itemize}
    \item $\sigma^{-1}(y) = \ln\left(\frac{y}{1-y}\right)$ is the logit function.
    \item $\delta = 10^{-5}$ is a numerical stabilizer for clipping.
    \item $\Phi_k^{(l)} \in \mathbb{R}^d$ and $\phi_k^{(l)} \in \mathbb{R}$ are learned task projection keys and biases defining the attractor basin orientation.
    \item $\lambda_l = \sigma(\lambda_{raw, l}) \in [0, 1]$ is the learned gating coefficient with a raw parameter $\lambda_{raw, l} \in \mathbb{R}$.
\end{itemize}

\paragraph{Gradient Stabilization Analysis:}
The G-CML formulation successfully tames chaotic gradient explosion. The derivative of the state $s^{(l)}$ with respect to the previous state $s^{(l-1)}$ contains a direct additive path of $(1 - \lambda_l)$:
\begin{equation}
\frac{\partial s^{(l)}}{\partial s^{(l-1)}} = (1 - \lambda_l)\mathbf{I} + \lambda_l \frac{\partial s_{cand}^{(l)}}{\partial s^{(l-1)}}.
\end{equation}
By initializing $\lambda_{raw, l}$ to a negative value (e.g., $-2.0$, yielding $\lambda_l \approx 0.12$), the network retains a strong skip connection of $1 - \lambda_l \approx 0.88$. This additive path guarantees a stable gradient flow directly through the layers, preventing the chaotic multiplier from compounding exponentially. This allows first-order optimizers like Adam to easily converge to robust parameter configurations.

\subsection{Task-Specific Routing and Weight Assembly}
In prior dynamic model-merging literature, sample-specific or task-specific coefficients are often averaged across heterogeneous batches to satisfy batch execution constraints. However, in highly sensitive chaotic systems, taking the arithmetic mean over a batch completely washes away task-specific trajectories, converting the dynamic system into a noisy static router.

We resolve this batch-averaging contradiction by introducing \textbf{Task-Specific Dynamic Routing}. Instead of performing batch-wide mean pooling over task-mixed features, we extract task-specific coefficients directly. For task $j$, the active lattice state $s_{k, j}^{(l)}$ maps directly to the task-specific merging coefficient:
\begin{equation}
\alpha_{k, j}(l) = R_k^{(l)} \cdot s_{k, j}^{(l)},
\end{equation}
where $R_k^{(l)} \in \mathbb{R}$ is a learned layer-wise scaling amplitude initialized to $0.3$.

The dynamically assembled expert weights for evaluating task $j$ at layer group $l$ are:
\begin{equation}
W_{merged, j}^{(l)} = W_{base}^{(l)} + \sum_{k=1}^K \alpha_{k, j}(l) V_k^{(l)}.
\end{equation}
This formulation preserves the exact dynamic, task-dependent nature of the chaotic state trajectories. When executing a multi-task workload, the backbone weights are dynamically assembled specifically for the target task features, resolving representational interference and preventing task signal dilution.

\paragraph{Practical Implementation via Task-Level Centroids:}
While the Coupled Map Lattice formulation theoretically supports sample-wise dynamic weight assembly (computing a unique $W_{merged, b}^{(l)}$ for each individual input sample $b$), doing so at inference time would require re-assembling and hot-swapping the entire network's parameters (e.g., 5.7M parameters) for each individual forward pass. This introduces massive memory-swapping and computational latency, which is highly impractical for resource-constrained edge computing. To resolve this execution bottleneck, in practice, we compute a single representative task-level feature centroid $\psi(x)_j$ (for example, by averaging the features over a tiny offline calibration set, or using a task-level running average centroid). We then use this centroid to dynamically merge and assemble the network weights exactly \emph{once per task/batch}. This task-level dynamic routing achieves an optimal balance, providing high representational customization for each task while maintaining standard execution speed during test-time.

Crucially, this centroid-based formulation does \emph{not} require an explicit Task ID at test-time, successfully maintaining a fully unsupervised and task-agnostic deployment. When an unlabeled batch of test samples arrives, we simply compute its feature centroid directly from the input features (i.e., $\bar{\psi} = \frac{1}{B} \sum_{b=1}^B \psi(x)_b$) in an unsupervised manner. If the incoming batch is heterogeneous (containing a mixture of tasks), practitioners can apply standard lightweight clustering (such as $K$-means) in the low-dimensional projected phase-space to automatically separate task clusters and compute centroids on-the-fly, completely bypassing any need for categorical Task IDs or oracle labels during inference.

Specifically, the computational cost of weight-space assembly is extremely small. In PyTorch, performing $W_{merged} = W_{base} + \sum \alpha_k V_k$ for a 5.7M-parameter model takes less than 2 milliseconds on a standard CPU and under 0.5 milliseconds on a modern GPU, as it consists of basic tensor additions and scalar multiplications. By performing this assembly once per task (or once per large batch of inputs), subsequent inference runs at the exact same high-throughput and zero-latency profile as a standard static model, bypassing the $100\times$ slower execution that sample-wise hot-swapping would introduce. This deliberate design choice guarantees that ChaosMerge delivers dynamic, customized multi-task representations with virtually zero test-time computational overhead.

\paragraph{Limitations of On-the-Fly Clustering in Heterogeneous Batches:}
While on-the-fly unsupervised clustering of input features (such as $K$-means in the sphere-projected 4-dimensional phase space) offers a theoretical pathway to execute ChaosMerge in completely task-agnostic, mixed-task batch environments without explicit Task IDs, we must explicitly acknowledge three key logical and practical bottlenecks:
\begin{itemize}
    \item \textbf{Cluster Count Estimation ($K$):} In truly task-agnostic deployments, the number of active tasks $K$ in an incoming batch is unknown. If the algorithm requires pre-specifying $K$ for clustering, it introduces a hidden hyperparameter dependency.
    \item \textbf{Inference Latency Multiplier:} If a mixed-task batch is split into $C$ clusters, the model must dynamically assemble and load $C$ distinct sets of weights, running separate forward passes for each partition. If $C$ is large, the execution throughput is divided by $C$, which can neutralize the benefits of standard batch execution.
    \item \textbf{Misclustering and Error Propagation:} Unsupervised clustering in a highly compressed 4-dimensional projected sphere space is highly prone to errors. Misclustering a sample results in evaluating it on weights merged specifically for a different task domain, which can cause catastrophic classification failure on that sample.
\end{itemize}

To evaluate these practical concerns empirically, we conducted a dedicated experiment using a heterogeneous test batch containing mixed-task samples from MNIST, FashionMNIST, CIFAR-10, and SVHN (128 samples per task; 512 samples total), and performed unsupervised spherical $K$-means ($K=4$) in G-CML's projected 4-dimensional phase space. Our results reveal significant bottlenecks:
\begin{enumerate}
    \item \textbf{Low Clustering Purity:} Due to the severe spatial overlap of the highly compressed 4-dimensional random projection space, the unsupervised spherical $K$-means clustering achieves a purity of only \textbf{45.31\%}.
    \item \textbf{Catastrophic Performance Degradation:} When evaluating each sample using the weights assembled for its assigned cluster, the downstream classification accuracy falls from the Oracle (perfect Task ID) baseline of \textbf{75.00\%} to only \textbf{45.31\%}. This severe \textbf{29.69\% absolute drop} in accuracy empirically proves that misclustering propagates catastrophically, as evaluating a sample on weights merged for an incorrect task domain results in near-zero accuracy on those samples.
    \item \textbf{Latency Multiplier:} Even on this small-scale test-batch, splitting the forward pass into $C=4$ sub-batches and running 4 distinct weight-assembly steps increases inference latency by \textbf{1.03$\times$} compared to a unified, single-assembly forward pass. This latency overhead would scale linearly with larger $C$ or larger backbones.
\end{enumerate}
These empirical findings demonstrate that unsupervised on-the-fly clustering is highly fragile. Consequently, while the centroid formulation is an exceptional optimization for task-specific deployments where task-level boundaries are known, running it task-agnostically via on-the-fly unsupervised clustering represents a major open research bottleneck. Addressing these challenges via robust cluster-number estimation (e.g., Dirichlet Process mixtures) and lightweight multi-centroid routing represents a vital direction for future dynamic deployment research.

\subsection{Parameter Complexity and Efficiency Analysis}
A major limitation of classical dynamic routing is parameter explosion, which frequently leads to the Overfitting-Optimizer Paradox when training on limited calibration data. ChaosMerge addresses this by utilizing a highly constrained, physically regularized CML structure. For a Vision Transformer backbone with $L=14$ layer groups and $K=4$ tasks, the trainable parameters are:
\begin{itemize}
    \item Initializer weight and bias: $W_{init} \in \mathbb{R}^{4 \times 4}$ and $b_{init} \in \mathbb{R}^4$ ($16 + 4 = 20$ parameters).
    \item Layer-wise coupling: $\gamma_{raw} \in \mathbb{R}^{14}$ ($14$ parameters).
    \item Scaling amplitudes: $R \in \mathbb{R}^{14 \times 4}$ ($56$ parameters).
    \item Attractor keys and biases: $\Phi \in \mathbb{R}^{14 \times 4 \times 4}$ and $\phi \in \mathbb{R}^{14 \times 4}$ ($224 + 56 = 280$ parameters).
    \item Layer-wise gating: $\lambda_{raw} \in \mathbb{R}^{14}$ ($14$ parameters).
\end{itemize}
This yields a total parameter count of \textbf{exactly 384 parameters}. This tiny footprint provides strong structural regularization, allowing ChaosMerge to be optimized extremely quickly on tiny calibration sets (such as 64 samples) without any risk of overfitting.
